\documentclass[10pt,letterpaper,final]{report}
\usepackage[margin=.75in]{geometry}
\usepackage[utf8]{inputenc}
\usepackage{amsmath}
\usepackage{amsfonts}
\usepackage{amssymb}
\usepackage{amsthm}
\renewcommand\qedsymbol{$\blacksquare$}
\usepackage{enumerate}
\usepackage{enumitem}
\usepackage{hyperref}
\usepackage{pdfpages}
\usepackage{graphics}
\usepackage{graphicx}
\usepackage{tikz}
\usepackage{tikz-qtree}
\usepackage{forest}
\usepackage{float}
\usetikzlibrary{automata,arrows}
\usetikzlibrary{shapes.multipart, positioning}
\usepackage{mdframed}
\usepackage[
  backend=biber,
  style=numeric-comp,
  sorting=none
]{biblatex}
\addbibresource{references.bib}

\usepackage{minted}
\usemintedstyle{algol_nu}
\usepackage{xltabular}
\usepackage{pdflscape}
\usepackage{tcolorbox}


% Based on template from COSC-417 at Towson University, originally authored by Marius Zimand

\begin{document}
\begin{center}
  \fbox{
    \vbox{
      \begin{flushleft}
        Jonathan Llovet \\  % authors' names
        EN.605.202.81 \\  %class
        2026-07-28 \\  % date
        Module 8 - Homework 8 \\ % ID
      \end{flushleft}
      \center{\Large{\textbf{Data Structures - Sorting}}}
      %\end{mdframed}
    } % end vbox
  } % end fbox
  \vline
\end{center}

\medskip

\noindent\textbf{Problem 1:}
\medskip

How many comparisons and interchanges (in terms of file size n) are performed by Simple insertion sort for the following files:
\begin{enumerate}
  \item A sorted file
  \item A file that is sorted in reverse order (that is, from largest to smallest)
  \item A file in which x[0], x[2], x[4]... are the smallest elements in sorted order, and in which x[1], x[3], x[5]... are the largest elements in sorted order, e.g.  [3, 14, 5, 15, 9, 18, 11, 19].
\end{enumerate}

\medskip
\noindent\textbf{Answer:}
\medskip

\noindent\textbf{Problem 2:}
\medskip

How many comparisons and interchanges (in terms of file size n) are performed by Shell Sort using increments 2 and 1 for the following files:
\begin{enumerate}
  \item A sorted file
  \item A file that is sorted in reverse order (that is, from largest to smallest)
  \item A file in which x[0], x[2], x[4]... are the smallest elements in sorted order, and in which x[1], x[3], x[5]... are the largest elements in sorted order, e.g. [3, 14, 5, 15, 9, 18, 11, 19].
\end{enumerate}

\medskip
\noindent\textbf{Answer:}
\medskip

\noindent\textbf{Problem 3:}
\medskip

Determine the number of comparisons (as a function of n and m) that are performed in merging two ordered files a and b of sizes n and m, respectively, by the merge method presented in the lecture, on each of the following sets of ordered files (a[i] refers the value in position i of file a, etc.):
\begin{enumerate}
  \item m=n and a[i] < b[i] < a[i+1], e.g. a=[6, 9, 12, 15, 29, 37] and b=[8, 10, 14, 25, 33, 45]
  \item m=n and a[n-1] < b[0], e.g. a=[2, 5, 9] and b=[12, 14, 16]
\end{enumerate}

\medskip
\noindent\textbf{Answer:}
\medskip

\noindent\textbf{Problem 4:}
\medskip

Determine the number of comparisons (as a function of n and m) that are performed in merging two ordered files a and b of sizes n and m, respectively, by the merge method presented in the lecture, on each of the following sets of ordered files:
\begin{enumerate}
  \item m=n and a[(n/2)-1] < b[0] < b[m-1] < a[(n/2)], e.g. a=[2, 5, 7,  55, 61, 72] and b=[9, 15, 17, 21, 29, 46]
  \item m=1 and b[0] < a[0]
  \item m=1 and a[n-1] < b[0], a[i] refers the value in position i of file a, etc.
\end{enumerate}

\medskip
\noindent\textbf{Answer:}
\medskip


\pagebreak

\section*{Source Code}
The full source code, including LaTeX, tooling, other artifacts, and git history is available on my Github repository for the class here:
\begin{center}
  \url{https://github.com/jllovet/jhu-en-605-202-su26-data-structures}
\end{center}

\printbibliography

\end{document}
